\section{Holy Orders}
\label{sec:orders}

Every Christian liturgy rests upon a quiet miracle of continuity: the Voice that spoke over bread and wine in the Upper Room still speaks across continents and centuries through mortal lips, and the Hands stretched out upon the Cross still bless, absolve, anoint, and shepherd through human hands laid upon a bent head. Before His Passion, Jesus Christ commanded the Twelve, ``Do this for a commemoration of me'' (Luke~22:19), and on the evening of His Resurrection He breathed the Holy Spirit upon them with His own mission from the Father (John~20:21--23). \textbf{Holy Orders} is the sacrament through which that apostolic mission continues to be exercised in the Church until the end of time, conferring upon baptized men chosen by the Church an indelible configuration to Christ in the three hierarchical degrees of the episcopate, the presbyterate, and the diaconate (CCC~1536; CIC~1983, cc.~1008--1009; CCEO~1990, c.~743).

\subsection{Scriptural Foundations: From Aaron and Melchisedech to the Laying on of Hands}
\label{subsec:orders-scripture}

The Epistle to the Hebrews sets forth the universal law of sacred mediation:
\begin{quotation}
For every high priest taken from among men is ordained for men in the things that appertain to God, that he may offer up gifts and sacrifices for sins\ldots\ Neither doth any man take the honour to himself, but he that is called by God, as Aaron was. So Christ also did not glorify himself, that he might be made a high priest: but he that said unto him: Thou art my Son: this day have I begotten thee. As he saith also in another place: Thou art a priest for ever, according to the order of Melchisedech. (Heb~5:1, 4--6)
\end{quotation}
Under the Old Covenant, God consecrated the house of Aaron to offer repeated sacrifices foreshadowing redemption (Exod~28--29; Lev~8), while Genesis had already placed before Sinai the mysterious figure of Melchisedech, king of Salem and ``priest of the most high God,'' who blessed Abraham ``bringing forth bread and wine'' (Gen~14:18--19). Jesus Christ fulfills the Levitical shadow and realizes the royal, bread-and-wine priesthood of Melchisedech. Because Levitical priests were mortal, death continually interrupted their office; Christ's priesthood is wholly other:
\begin{quotation}
But this, for that he continueth for ever, hath an everlasting priesthood: Whereby he is able also to save for ever them that come to God by him; always living to make intercession for us. For it was fitting that we should have such a high priest, holy, innocent, undefiled, separated from sinners, and made higher than the heavens: Who needeth not daily (as the other priests) to offer sacrifices, first for his own sins, and then for the people's: for this he did once, in offering himself. (Heb~7:24--27)
\end{quotation}
The uniqueness of Christ's priesthood does not abolish ministerial priesthood in the Church; it defines it. An ordained priest neither replaces an absent Christ nor adds a second sacrifice to Calvary. Because the High Priest lives forever, He acts sacramentally through ordained men whom He makes living instruments of His single, eternal mediation.

In the Gospels, Christ prepares and institutes this ministry through deliberate acts restricted to the Twelve. He calls them after a night of prayer to be sent with His authority (Mark~3:13--15; Luke~6:13). At the Last Supper He delivers the Eucharistic mandate: ``This is my body, which is given for you. Do this for a commemoration of me'' (Luke~22:19; 1~Cor~11:24--25). On Easter evening, as seen in the chapter on Penance (Section~\ref{sec:penance}), the risen Lord joins apostolic mission directly to the gift of the Spirit, sending the Eleven as the Father sent Him, breathing the Holy Ghost upon them, and entrusting to them the authority to forgive or retain sins (John~20:21--23). Before ascending, He commissions them to teach and baptize all nations (Matt~28:18--20) and entrusts to Peter and the apostolic college the shepherding of His flock (John~21:15--17; Matt~16:18--19; 18:18).

Immediately after Pentecost, Acts and the Pastoral Epistles show this apostolic office transmitted through a definite visible rite: prayer joined to the \textbf{imposition of hands} (\latin{impositio manuum}; Greek \latin{cheirotonia}). When the Jerusalem community grew, the Twelve had the brethren choose seven men of good reputation, full of the Holy Spirit and wisdom: ``These they set before the apostles: and they praying, imposed hands upon them'' (Acts~6:6)---the apostolic origin of the diaconate in charitable, liturgical, and kerygmatic service. At Antioch, the prophets and teachers fast, pray, and impose hands upon Barnabas and Saul at the Spirit's command (Acts~13:2--3). Returning through Asia Minor, Paul and Barnabas establish stable local governance: ``And when they had ordained to them priests in every church and had prayed with fasting, they commended them to the Lord, in whom they believed'' (Acts~14:22, Douay-Rheims; Greek~14:23). Paul charges Titus in Crete to ``ordain priests in every city, as I also appointed thee'' (Titus~1:5). Above all, the letters to Timothy establish that this imposition of hands confers an abiding interior gift (\latin{charisma}, Vulgate \latin{gratia}): ``Neglect not the grace that is in thee, which was given thee by prophecy, with imposition of the hands of the priesthood'' (1~Tim~4:14), and again, ``I admonish thee that thou stir up the grace of God which is in thee by the imposition of my hands'' (2~Tim~1:6).

\subsection{Patristic Tradition and the Historical Clarification of the Essential Rite}
\label{subsec:orders-tradition}

At the dawn of the second century, while led in chains from Antioch to martyrdom in Rome, St.~Ignatius of Antioch testifies in his authentic Shorter Recension that the threefold hierarchy of bishop, presbyters, and deacons belongs to the visible constitution of the Church:
\begin{quotation}
In like manner, let all reverence the deacons as an appointment of Jesus Christ, and the bishop as Jesus Christ, who is the Son of the Father, and the presbyters as the sanhedrim of God, and assembly of the apostles. Apart from these, there is no Church. (\emph{Epistle to the Trallians}~3)
\end{quotation}
Deacons, he reminds the Trallians, are not mere stewards of food and drink, ``but servants of the Church of God'' and ministers of Christ's mysteries (\emph{Trallians}~2). To the Smyrnaeans, Ignatius joins sacramental communion directly to the bishop:
\begin{quotation}
See that ye all follow the bishop, even as Jesus Christ does the Father, and the presbytery as ye would the apostles; and reverence the deacons, as being the institution of God. Let no man do anything connected with the Church without the bishop. Let that be deemed a proper Eucharist, which is [administered] either by the bishop, or by one to whom he has entrusted it. Wherever the bishop shall appear, there let the multitude [of the people] also be; even as, wherever Jesus Christ is, there is the Catholic Church. (\emph{Epistle to the Smyrnaeans}~8)
\end{quotation}

In the fourth century, St.~John Chrysostom contemplates the dignity of the Christian priesthood in his six books \emph{On the Priesthood} (\emph{De sacerdotio}~III), marveling that mortal men clothed in flesh are entrusted with a heavenly office---invoking the Holy Spirit at the altar and exercising a power to bind and loose sins granted neither to angels nor to archangels. Preaching on 2~Tim~1:6, Chrysostom teaches that the Spirit's gift imparted through the laying on of hands endures within the minister and demands unflagging zeal:
\begin{quotation}
As fire requires fuel, so grace requires our alacrity, that it may be ever fervent. ``I put thee in remembrance that thou stir up the gift of God, that is in thee by the putting on of my hands,'' that is, the grace of the Spirit, which thou hast received, for presiding over the Church, for the working of miracles, and for every service. (\emph{Homilies on 2~Timothy}~I)
\end{quotation}
Commenting on 1~Tim~4:14, he further preserves the ancient canonical distinction between bishop and presbyter: ``He speaks not here of Presbyters, but of Bishops. For Presbyters cannot be supposed to have ordained a Bishop'' (\emph{Homilies on 1~Timothy}~XIII). At the close of the patristic age, Pseudo-Dionysius the Areopagite (\emph{De ecclesiastica hierarchia}~V) contemplated the triad of hierarchs (bishops), priests (presbyters), and ministers (deacons) as the earthly icon of the celestial hierarchy, noting explicitly that three perfecting rites belong exclusively to the bishop: the ordination of all sacred ministers, the consecration of holy chrism (\latin{myron}), and the consecration of the altar (PG~3, cols.~504--508).

Medieval scholastic theology---systematized in St.~Thomas Aquinas (\emph{ST}~III, qq.~63, 65) and the \emph{Supplementum} (qq.~34--40, drawn from his \emph{Scriptum super Sententiis}~IV)---ordered every degree of Holy Orders to the Holy Eucharist, either directly to consecrating Christ's true Body and Blood or ministerially to serving that sacrifice and preparing the Mystical Body. Because medieval Latin ordination rites incorporated symbolic handings-over of liturgical vessels and books (\latin{traditio} or \latin{porrectio instrumentorum}---paten with bread and chalice with wine for presbyters, Book of the Gospels for deacons), Latin scholastics and the Council of Florence's \emph{Decree for the Armenians} (\emph{Exsultate Deo}, 1439, \S15) highlighted those ceremonies alongside the ancient imposition of hands, whereas the Eastern Churches had always conferred Orders solely through imposition of hands (\latin{cheirotonia}) and consecratory prayer.

Against sixteenth-century denials of a visible sacrificing priesthood, the \textbf{Council of Trent} (Session~XXIII, 15~July 1563) defined the permanent Catholic doctrine: in the New Testament there is a visible and external priesthood instituted by Christ with power to consecrate and offer His true Body and Blood and to forgive and retain sins (ch.~1, can.~1); Holy Orders is truly a sacrament conferring grace and imprinting an indelible character so that an ordained priest never becomes a layman again (ch.~3--4, cans.~3--4); and a divinely instituted hierarchy exists comprising bishops, presbyters, and ministers, in which bishops are superior to presbyters and hold the proper power to confirm and ordain (ch.~4, cans.~6--7). The \emph{Roman Catechism} (1566, Part~II, \emph{On Orders}) distinguishes the twofold authority entrusted to the Church's sacred hierarchy: the \emph{power of order}, directed to Christ's true Body in the Holy Eucharist, and the \emph{power of jurisdiction}, directed to governing Christ's Mystical Body and guiding the faithful unto salvation.

Three modern magisterial landmarks brought the theology of the essential rite and the three degrees to complete clarity.

First, in \emph{Apostolicae Curae} (13~September 1896), Pope Leo~XIII examined the Edwardine Anglican Ordinal and stated the governing law of sacramental signification and intention:
\begin{quotation}
All know that the sacraments of the New Law, as sensible and efficient signs of invisible grace, ought both to signify the grace which they effect, and effect the grace which they signify. Although the signification ought to be found in the whole essential rite---that is to say, in the matter and form---it still pertains chiefly to the form; since the matter is the part which is not determined by itself, but which is determined by the form. And this appears still more clearly in the Sacrament of Orders, the matter of which, in so far as We have to consider it in this case, is the imposition of hands, which indeed by itself signifies nothing definite, and is equally used for several Orders and for Confirmation. (\emph{Apostolicae Curae})
\end{quotation}
Because the Edwardine Ordinal deliberately removed every mention of sacrifice and sacrificing priesthood, its words failed to signify definitely either the presbyterate or the high priesthood of the episcopate; and that deliberate alteration of the Catholic rite manifested a defect of intention:
\begin{quotation}
When any one has rightly and seriously made use of the due form and the matter requisite for effecting or conferring the sacrament he is considered by the very fact to do what the Church does\ldots\ On the other hand, if the rite be changed, with the manifest intention of introducing another rite not approved by the Church and of rejecting what the Church does, and what by the institution of Christ belongs to the nature of the sacrament, then it is clear that not only is the necessary intention wanting to the sacrament, but that the intention is adverse to and destructive of the sacrament. (\emph{Apostolicae Curae})
\end{quotation}
On this double defect of form and intention, Leo~XIII declared ordinations conferred according to the Anglican rite absolutely null and utterly void.

Second, Pope Pius~XII's Apostolic Constitution \emph{Sacramentum Ordinis} (30~November 1947), followed by St.~Paul~VI's Apostolic Constitution \emph{Pontificalis Romani} (18~June 1968), settled the medieval Latin debate over the \latin{traditio instrumentorum}. Pius~XII declared that for the diaconate, presbyterate, and episcopate in the Latin Church, the \textbf{sole essential matter} is the bishop's imposition of hands, and the \textbf{sole essential form} consists in the specified consecratory prayer determining the grace and office of that Order. Latin and Eastern sacramental theology thus stand in full unity on the apostolic core of the rite, while anointing, vesting, and the presentation of vessels or Gospels remain rich explanatory ceremonies.

Third, the \textbf{Second Vatican Council} and \textbf{Pope Benedict~XVI} clarified the exact relation between the common priesthood, the two priestly degrees, and the diaconate. In \emph{Lumen Gentium}~10, the Council teaches that the common priesthood of all the baptized and the ministerial priesthood differ \textbf{in essence and not only in degree} (\latin{essentia et non gradu tantum}), yet are ordered one to another as distinct participations in the one priesthood of Christ (\emph{Presbyterorum Ordinis}~2, 5). In \emph{Lumen Gentium}~21, the Council resolved medieval hesitation by teaching that \textbf{episcopal consecration confers the fullness of the sacrament of Orders} (\latin{plenitudo sacramenti Ordinis}), the high priesthood and summit of sacred ministry, conferring the offices of sanctifying, teaching, and governing to be exercised in hierarchical communion with the Head and members of the episcopal college. In \emph{Lumen Gentium}~29, the Council restored the \textbf{permanent diaconate} in the Latin Church, noting that upon deacons hands are imposed ``not unto the priesthood, but unto a ministry of service'' (\latin{non ad sacerdotium, sed ad ministerium}). To align Latin canon law strictly with that doctrine and CCC~875 and~1554, Pope Benedict~XVI's motu proprio \emph{Omnium in Mentem} (26~October 2009) amended CIC~1983 canons~1008 and~1009 (adding c.~1009~\S3): bishops and presbyters receive the mission and capacity to act \latin{in persona Christi Capitis} (in the person of Christ the Head), whereas deacons are empowered to serve the People of God in the \latin{diakonia} (ministry) of liturgy, word, and charity. The diaconate is truly a sacramental degree of Holy Orders, yet strictly \textbf{non-priestly}: a deacon cannot consecrate the Eucharist, absolve sins in Penance, or administer the Anointing of the Sick.

\subsection{Sacramental Structure: Matter, Form, Minister, Subject, and Discipline}
\label{subsec:orders-structure}

Every valid ordination rests upon four determinate elements fixed by Christ and authenticated by the Church.

\textbf{Essential Matter.} The sole essential matter for all three degrees is the consecrating bishop's silent \textbf{imposition of hands} upon the head of the ordinand (\emph{Sacramentum Ordinis}; \emph{Pontificalis Romani}; CIC~1983, c.~1009~\S2; CCEO~1990, c.~744). In episcopal ordination, the principal consecrator and co-consecrating bishops impose hands silently; in presbyteral ordination, attending presbyters lay hands on the candidate after the bishop to signify incorporation into the presbyterium, though the bishop's imposition and prayer alone effect the sacrament.

\textbf{Essential Form.} The sole essential form consists in the determinate \textbf{consecratory prayer} prescribed in the approved liturgical books of each Church \latin{sui iuris}---in the Roman Rite, the central petition within the Prayer of Ordination for bishop, presbyter, or deacon (\emph{Pontificalis Romani}, 1968), and in Eastern Catholic rites, the corresponding epicletic prayer (\latin{cheirotonia}) invoking divine grace to elevate the candidate to the named Order. Because imposition of hands is common to all three degrees, the consecratory prayer strictly specifies the Order conferred.

\textbf{Minister.} A \textbf{validly consecrated bishop alone} is the minister of every degree of Holy Orders (Trent, Sess.~XXIII, can.~7; CIC~1983, c.~1012; CCEO~1990, c.~744). No presbyter can validly confer any degree of Orders. For liceity in episcopal ordination, canon law strictly requires a pontifical mandate (or prescribed patriarchal/synodal election and assent in Eastern law) and ordinarily requires at least two co-consecrating bishops to manifest collegial communion (CIC~1983, cc.~1013--1014; CCEO~1990, cc.~745--746). Consecrating a bishop without pontifical mandate, while valid if the essential rite and intention are preserved, incurs automatic excommunication reserved to the Apostolic See (CIC~1983, c.~1387).

\textbf{Subject.} Only a \textbf{validly baptized male} validly receives sacred ordination (CIC~1983, c.~1024; CCEO~1990, c.~754). In \emph{Ordinatio Sacerdotalis} (22~May 1994), St.~John~Paul~II declared definitively that the Church has no authority whatsoever to confer priestly ordination on women and that this judgment must be definitively held by all the faithful. This reservation rests not on any inequality of personal dignity or holiness---the Blessed Virgin Mary, who received neither apostolic office nor ministerial priesthood, stands above every apostle and angel---but on fidelity to Christ's own free choice of the Twelve, unbroken apostolic practice, and the sacramental icon of Christ the Bridegroom and Head (\latin{in persona Christi Capitis}). For licit ordination, canon law requires free vocation attested in writing, sound faith, proven integrity, proper age and formation, and freedom from irregularities and impediments (CIC~1983, cc.~1025--1050; CCEO~1990, cc.~758--762).

\begin{fourstable}{Degree of Order}{Configuration \& Rank}{Essential Rite \& Minister}{Proper Ministry \& Discipline}
Episcopate & Fullness of Orders (\latin{summum sacerdotium}); acts \latin{in persona Christi Capitis} & Bishop's silent imposition of hands + episcopal consecratory prayer; valid bishop (ordinarily with two co-consecrators) & Sanctifies, teaches, and governs a particular Church in collegial communion; sole minister of Holy Orders; exclusively celibate in East and West\\
Presbyterate & Co-worker of the episcopal order; acts \latin{in persona Christi Capitis} & Bishop's silent imposition of hands + presbyteral consecratory prayer; valid bishop alone & Consecrates Eucharist, absolves in Penance, administers Anointing, preaches, shepherds flock; celibate in Latin, Syro-Malabar, Syro-Malankara, and Alexandrian Churches; married men admissible in most other Eastern Churches\\
Diaconate & Non-priestly sacred Order (\latin{non ad sacerdotium}); acts in \latin{diakonia} & Bishop's silent imposition of hands + diaconal consecratory prayer; valid bishop alone & Assists at altar, proclaims Gospel, preaches, baptizes, assists at Latin marriages, leads works of charity; open to mature married men (permanent) or celibate men\\
\end{fourstable}

\textbf{Instituted Lay Ministries Distinct from Orders.} By his motu proprio \emph{Ministeria quaedam} (15~August 1972), St.~Paul~VI suppressed clerical tonsure, minor orders, and the subdiaconate in the Latin Church, bound entry into the clerical state strictly to \textbf{ordination to the diaconate}, and retained lector and acolyte as \textbf{instituted lay ministries} grounded in Baptism. Pope Francis's motu proprio \emph{Spiritus Domini} (10~January 2021) amended CIC~1983 canon~230~\S1 so that lay women as well as lay men may receive stable institution as lectors and acolytes (joined by catechist under \emph{Antiquum ministerium}, 2021). Because these ministries flow from the common priesthood of Baptism, opening them to all the lay faithful sharpens the boundary separating lay ministry from ordained Holy Orders.

\textbf{Clerical Celibacy Across Latin and Eastern Catholic Discipline.} Perfect and perpetual continence is not required by the very nature of the priesthood, as the primitive Church and Eastern Tradition prove (\emph{Presbyterorum Ordinis}~16; Paul~VI, \emph{Sacerdotalis caelibatus}~14, 17). Yet celibacy chosen for the kingdom of heaven bears deep Christological, ecclesiological, and eschatological fittingness, conforming the minister to Christ's virginal self-gift, freeing him for undivided spousal charity toward the Church, and pointing visibly to the age to come (\emph{Sacerdotalis caelibatus}~19--34; John~Paul~II, \emph{Pastores dabo vobis}~29).

In the \textbf{Latin Church}, ecclesiastical law requires celibacy of all candidates for the presbyterate and of unmarried candidates for the permanent diaconate (CIC~1983, cc.~277~\S1, 1037, 1042~1\textdegree), with narrow papal dispensations for married former ministers of separated communions (\emph{Sacerdotalis caelibatus}~42; Benedict~XVI, \emph{Anglicanorum coetibus}). Married men of at least thirty-five years, with their wife's written consent, may be ordained permanent deacons (CIC~1983, cc.~1031~\S2, 1050~3\textdegree; \emph{Lumen Gentium}~29).

In the \textbf{Eastern Catholic Churches}, CCEO~1990 canon~373 honors both clerical celibacy and the venerable tradition of married clergy, exhorting married clerics to exemplary chastity and family life (CCEO~1990, cc.~374--375). Particular law of each Church \latin{sui iuris} governs the admission of married men to sacred orders (CCEO~1990, cc.~758~\S3, 762). While most Byzantine, Antiochene, Armenian, and Chaldean Churches admit married men to the diaconate and presbyterate, mandatory priestly celibacy remains the general particular-law discipline not only in the Syro-Malabar and Syro-Malankara Churches, but also in the Alexandrian Catholic Churches (Coptic, Ethiopian, and Eritrean, subject to rare hierarchical dispensation). Governing diaspora practice outside traditional Eastern territories, the Congregation for the Eastern Churches' decree \emph{Pontificia Praecepta de clero uxorato orientali} (14~June 2014, approved by Pope Francis on 23~December 2013, AAS~106 [2014], 496--499) restored Eastern Hierarchs' faculty to ordain married men in those territories for Churches possessing that tradition after informing the local Latin bishop in writing.

Two universal rules unite East and West without exception: only celibate priests are consecrated bishops (CCEO~1990, c.~180; \emph{Sacerdotalis caelibatus}~40), and sacred ordination creates a \textbf{diriment impediment to subsequent marriage} (Trent, Sess.~XXIV, can.~9; CIC~1983, c.~1087; CCEO~1990, c.~804). No cleric in sacred orders may validly marry or remarry after ordination absent a rare papal dispensation (CIC~1983, cc.~291, 1078~\S2; CCEO~1990, c.~795).

\subsection{Sacramental Reality, Indelible Character, Grace, and Order of Ends}
\label{subsec:orders-effects}

\textbf{\latin{Res et sacramentum} and Permanent Formal Change.} The immediate, objective spiritual reality effected by every valid ordination is the \textbf{indelible sacramental character} of deacon, presbyter, or bishop (Trent, Sess.~VII, can.~9; Sess.~XXIII, can.~4; CCC~1582). Metaphysically, this character is a \textbf{permanent accidental formal elevation} of the intellective soul---an instrumental spiritual power configuring the recipient to Christ the Eternal Priest and Servant (\emph{ST}~III, q.~63, aa.~2--5). St.~Thomas explains its imperishability:
\begin{quotation}
Now Christ's Priesthood is eternal, according to Ps.~109:4: ``Thou art a priest for ever, according to the order of Melchisedech.'' Consequently, every sanctification wrought by His Priesthood, is perpetual, enduring as long as the thing sanctified endures\ldots\ Since, therefore, the subject of a character is the soul as to its intellective part, where faith resides\ldots\ it is clear that, the intellect being perpetual and incorruptible, a character cannot be blotted out from the soul. (\emph{ST}~III, q.~63, a.~5)
\end{quotation}
Consequently, \textbf{no degree of Holy Orders can ever be repeated}, and \textbf{valid ordination never becomes invalid} (CIC~1983, c.~290; CCEO~1990, c.~394). Even a priest juridically dismissed from the clerical state retains his priestly character; in urgent danger of death he validly and licitly absolves any penitent from all sins and censures (CIC~1983, cc.~292, 976).

\textbf{\latin{Res tantum} and Proper Sacramental Grace.} In a duly disposed recipient free of mortal sin, Holy Orders confers an increase of sanctifying grace joined to the \textbf{special sacramental grace of office} (\latin{gratia status}, CCC~1585--1589; 2~Tim~1:6): strength for the bishop to teach, govern, and lay down his life; pastoral charity and reverence at the altar for the presbyter; humility, fidelity to the Gospel, and compassionate service of the poor for the deacon. If an obstacle (\latin{obex}) initially blocks fruitfulness, the character is still validly imprinted, and grace revives upon repentance.

\textbf{Order of Ends.} The teleology of Holy Orders unfolds across four ordered horizons:
\begin{itemize}
  \item \textbf{Primary proper end:} perpetuation of Christ's apostolic ministry in the Church---consecrating and offering the Eucharistic sacrifice, administering the sacraments, proclaiming the Word with authority, shepherding the flock, and serving in sacred \latin{diakonia}.
  \item \textbf{Intrinsic secondary ends:} personal sanctification of the minister through pastoral charity (\emph{Presbyterorum Ordinis}~13--14), visible unity of each particular Church around its bishop, and defense of the community against heresy and schism (\emph{ST}~III, q.~65, a.~1).
  \item \textbf{Contingent ends:} historically adapted pastoral structures, missionary outreach, and charitable administration suited to each age and culture.
  \item \textbf{Ultimate end:} building up the Mystical Body of Christ unto the measure of His fullness and eternal praise in the heavenly sanctuary (Eph~4:12--13; Apoc~5:9--10).
\end{itemize}

\begin{threestable}{Structural Dimension}{Doctrinal Substance}{Canonical and Liturgical Specification}
Essential matter & Silent imposition of hands by the consecrating bishop upon the ordinand's head & Fixed by Pius~XII (\emph{Sacramentum Ordinis}, 1947) and Paul~VI (\emph{Pontificalis Romani}, 1968); explanatory rites unfold signification without constituting validity\\
Essential form & Specific consecratory prayer invoking the Holy Spirit for the exact degree conferred & Distinct essential petitions for episcopate, presbyterate, and diaconate in approved Latin and Eastern books (CIC~c.~1009~\S2; CCEO~c.~744)\\
Minister & Validly consecrated bishop alone & Sole valid minister of all three degrees (CIC~c.~1012; CCEO~c.~744); episcopal consecration requires pontifical mandate and ordinarily two co-consecrators (CIC~cc.~1013--1014)\\
Subject & Validly baptized male alone, freely called and tested by the Church & Divine reservation of priestly ordination to men (\emph{Ordinatio Sacerdotalis}; CIC~c.~1024; CCEO~c.~754); celibacy discipline varies between Latin and most Eastern Churches\\
\latin{Res et sacramentum} & Indelible diaconal, presbyteral, or episcopal character & Permanent accidental formal elevation of the soul (\emph{ST}~III, q.~63, a.~5); never repeated, never invalidated upon loss of clerical state (CIC~c.~290; CCEO~c.~394)\\
\latin{Res tantum} & Increase of sanctifying grace plus proper sacramental grace of office & Conforms bishops and presbyters to Christ the Head (\latin{in persona Christi Capitis}) and deacons to Christ the Servant (\latin{diakonia}; \emph{Omnium in Mentem}, 2009; CIC~c.~1009~\S3)\\
\end{threestable}
